\chapter{نظام الملفات}
\label{chap:fs}

يتولى \indextext{نظام الملفات (File System)} تنظيم البيانات وحفظها بصفة دائمة على وسائل التخزين كالأقراص الصلبة.
ينظم نظام الملفات في \texttt{xv6} البيانات بداخل سبع طبقات متكاملة:
\begin{enumerate}
\item \textbf{طبقة القرص (Disk Layer)}: القراءة والكتابة المباشرة بكتل القرص عبر واجهة العتاد.
\item \textbf{طبقة المخزن المؤقت (Buffer Cache)}: التخزين المؤقت لكتل القرص الأكثر استمالاً في الذاكرة لتسريع الوصول وتحديد التزامن.
\item \textbf{طبقة سجل العمليات (Logging Layer)}: ضمان استعادة اتساق نظام الملفات عند انقطاع التيار الكهربائي أو الفشل المفاجئ.
\item \textbf{طبقة العقد الفهرسية (Inode Layer)}: إدارة العقد الفهرسية وتوزيع كتل البيانات الخاصة بكل ملف.
\item \textbf{طبقة المجلدات (Directory Layer)}: ربط الأسماء النصية بأرقام العقد الفهرسية.
\item \textbf{طبقة المسارات (Pathname Layer)}: تصفح المجلدات الشجرية مثل \texttt{/a/b/c} وتحليل المسارات.
\item \textbf{طبقة واصفات الملفات (File Descriptor Layer)}: تقديم واجهة الموارد الموحدة للعمليات والتطبيقات.
\end{enumerate}
